\chapter{Introduction}
\label{ch:introduction}

Password authentication remains the dominant credential mechanism for
consumer web services, in spite of decades of proposals for alternatives.
Its persistence is a consequence of deployability, familiarity, and a near-zero
marginal cost. Its cost is a long history of database breaches that
repeatedly expose hundreds of millions of hashed credentials to offline
guessing attacks. Once a hash file has left the server, the guessing rate
is bounded only by the attacker's compute, and any weakness in the
underlying password distribution translates directly into recovered
plaintexts.

\textcite{juels2013honeywords} proposed \emph{honeywords} as a detection
layer over this failure mode. For each account the authentication server
stores a \emph{sweetword set} of $k$ candidate hashes, one of which is the
real password and the other $k-1$ of which are decoys, together with a
separate index held by a hardened \emph{HoneyChecker} that identifies
which position is real. A successful login with a decoy is a signal that
the hash file has been read offline and the attacker has begun trying
candidates. The scheme does not prevent guessing; it converts silent
database compromise into an event that the server can detect and act on.

The security of the scheme rests on a property called
\emph{flatness}: given the sweetword set, an attacker who has read the
hash file must not be able to identify the real password with probability
much better than $1/k$. Under perfect flatness, all $k$ candidates are
indistinguishable to the attacker, and every login attempt against a decoy
raises the alarm with the same probability. Flatness is what turns
honeywords from a labeling exercise into a detection guarantee.

\section{Motivation}

Flatness, in most of the honeyword literature, is measured against
attackers that no longer represent the state of the art. Traditional
evaluations rank sweetwords using frequency counts, probabilistic
context-free grammars, or Markov models, and honeywords are declared
flat when the real password does not stand out under those rankers.
\textcite{dani2024passfilter} showed that a modestly sized
character-level convolutional neural network trained as a binary
classifier defeats this notion on the honeyword schemes it was tested
against, including the representation-learning generators of
\textcite{dionysiou2021honeygen}. Under a learned discriminator, the
attacker success rate is substantially higher than $1/k$, sometimes
by an order of magnitude, and sweetwords that pass heuristic evaluation
turn out to be individually plausible but collectively distinguishable.

This gap has two consequences that motivate the rest of the thesis.
First, the current honeyword generators have almost never been benchmarked
against a neural attacker, so it is not known how badly they fail once
one is admitted into the threat model. Second, no construction in the
literature restores flatness against such an attacker without redesigning
the generator itself, which would leave existing deployments unprotected.
Both issues need to be resolved for honeyword-based breach detection
to remain a credible defense.

\section{Approach}

We take the same discriminator that \textcite{dani2024passfilter}
introduced as an attack tool and reuse it, unmodified in architecture,
as a \emph{defensive oracle} during sweetword construction. A candidate
sweetword set is produced by any of the standard mechanisms, scored by
the oracle, and either accepted or discarded depending on whether the
induced score distribution over the set satisfies a small collection of
measurable flatness conditions. Sets that fail are regenerated up to a
bounded budget. The classifier that would otherwise identify the real
password thus becomes the gatekeeper that prevents any set with an
identifiable real password from being stored.

The construction is deliberately generator-agnostic. We apply it as a
post-processing layer over ten honeyword generation methods, spanning
rule-based tweaking, representation-learning approaches derived from
\textcite{dionysiou2021honeygen}, and a self-trained GAN. Each method
is trained once at a fixed training-set size and evaluated on a common
test set of sweetword sets, so that comparisons across generators are
made on the same footing. The primary evaluation uses a CNN attacker
trained on HoneyGen; robustness is then checked against two
further attacker classes -- a mixed-corpus CNN trained on several
decoy families and a frequency ranker -- which together distinguish
genuine flatness from single-attacker artifacts.

The primary security metric is the attacker's success rate at guess
budget five, $\mathrm{ASR}@5$: the fraction of accounts an attacker
compromises when it is allowed five guesses among the $k = 20$
sweetwords of an account. We report the full curve $\mathrm{ASR}@t$ for
$t = 1, \dots, 20$, but summarise on $t = 5$ because it matches the
honeyword threat model. After cracking a stolen hash file offline, an
attacker still has to submit a recovered sweetword to the live service
to use it; submitting a honeyword trips the honeychecker, and most
online services additionally lock an account after a small number of
failed attempts, with five a common threshold
\parencite{dani2024passfilter}. $\mathrm{ASR}@5$ therefore measures what
an attacker can achieve within one realistic lockout window. The
extremes are less informative as a single headline: $\mathrm{ASR}@1$
reflects only the single best guess, while $\mathrm{ASR}@{20}$ equals
$1.0$ by construction once all $k = 20$ sweetwords have been tried. For
a perfectly flat set the attacker does no better than random, giving
$\mathrm{ASR}@5 = 5/20 = 0.25$.

\section{Contributions}

The thesis makes the following contributions.

\begin{enumerate}

    \item \textbf{Attacker-relative set-level flatness.}
    We define flatness as a property of the probability distribution
    that a learned attacker induces over a sweetword set, and we
    characterize it by four quantities computed from that distribution:
    rank-based flatness, normalized entropy, top and real gaps, and an
    $\varepsilon$-flatness bound on the largest renormalized score in
    the set. These replace the frequency-based flatness notions that
    dominate the earlier literature and make the property operational
    under any binary discriminator.

    \item \textbf{Hybrid flatness filtering.}
    We propose a two-stage rejection-sampling algorithm that combines
    candidate-level pre-screening with set-level acceptance.
    Pre-screening removes individually distinguishable decoys before a
    set is assembled; set-level acceptance discards sets whose induced
    distribution violates any of the flatness conditions. The algorithm
    operates as a post-processing layer and does not require modifying
    the underlying generator or retraining the attacker.

    \item \textbf{Benchmark of ten honeyword generators against a neural
    attacker.} We evaluate Chaffing-by-Tweaking, Chaffing-with-Password-Model,
    Chaffing-with-Hybrid-Model, HoneyGen Baseline, HoneyGen Set-Filter,
    HoneyGen Prescreen-Add, HoneyGen Prescreen-Ratio, HoneyGen Hybrid-Add,
    HoneyGen Hybrid-Ratio, and a self-trained GAN, on the RockYou corpus
    with $k=20$. Reporting the same attacker success rate under a common
    experimental protocol lets us separate genuine flatness gains from
    distributional artifacts of the generator, an effect that turns out
    to matter for the GAN and that we discuss in \Cref{ch:results}.

    \item \textbf{A defensive reinterpretation of the neural-attacker
    paradigm.} The thesis shows that a discriminator of the type used
    by \textcite{dani2024passfilter} as an attack can be repurposed,
    without architectural change, into the enforcement mechanism of a
    flatness guarantee. What defeats current honeyword schemes at test
    time becomes what protects them at storage time.

\end{enumerate}

\section{Scope and Assumptions}

The guarantees this thesis provides are model-relative. Filtering is
performed against one specific attacker, the character-level CNN
described in \Cref{ch:method}, and the flatness conditions constrain
that attacker's induced distribution over the sweetword set.
A different attacker, with a different architecture or with access to
side information beyond the sweetword strings, may recover
distinguishability. \Cref{ch:conclusion} returns to this point when
discussing the limits of the framework.

The evaluation is restricted to leaked passwords from the RockYou
corpus and to sweetword sets of size $k=20$. Only the password string
itself is used as input by the generators and the attackers.
\Cref{ch:results} reports how far the conclusions carry over across
attacker classes, from the HoneyGen-tuned CNN to a mixed-corpus CNN
and a frequency ranker; extending the benchmark to a second corpus
such as \emph{darkc0de} is left as future work
(\Cref{ch:conclusion}).

The remainder of the thesis proceeds as follows. \Cref{ch:background}
sets out the honeyword primitive, the HoneyChecker architecture, and
the character-level convolutional networks that both the attack and the
defense rely on. \Cref{ch:related} reviews the honeyword generators
we benchmark against and situates \textcite{dani2024passfilter} within
the wider line of neural password modeling. \Cref{ch:method} develops
the flatness formalism, the CNN training procedure, the ten generators,
the two-stage filtering algorithm, and the evaluation protocol.
\Cref{ch:results} reports the ASR curves across
the ten generators, isolates the effect of filtering on both security
and the flatness criteria, weighs it against the acceptance rate as a
deployability requirement, checks robustness across attacker classes,
and discusses a distributional artifact in the GAN results that limits
how far its low ASR should be trusted. \Cref{ch:conclusion} draws the
argument together and lists the open questions the framework leaves
behind.
